Для изображений химических структур в программном обеспечении используется тот же общий принцип адаптивной предобработки, однако с рядом изменений, связанных со спецификой задачи OCSR. Если в случае обычного текста основная цель предобработки состоит в повышении читаемости символов для OCR-движка, то при обработке химических формул и молекулярных схем ключевым требованием становится сохранение топологии структуры: взаимного расположения атомов, связей, индексов, зарядов и других графических элементов. По этой причине в модуле предобработки химических структур применяется та же последовательность этапов, что и для текстовых изображений, но параметры обработки выбираются другим способом.

Главное отличие состоит в том, что адаптация параметров ориентируется не на характерный размер текстового символа, а на характерный масштаб химического объекта на изображении. Это позволяет подстраивать увеличение, бинаризацию и морфологическую обработку под размер молекулярной схемы, а не под особенности печатного текста. Масштабирование по-прежнему используется для выравнивания качества входных данных, однако его задача состоит не столько в повышении разборчивости отдельных символов, сколько в приведении тонких связей и мелких обозначений к диапазону размеров, более удобному для последующего распознавания.

Дополнительные отличия касаются интенсивности локальных преобразований. Усиление контраста, шумоподавление и повышение резкости в режиме OCSR применяются в более консервативной форме, поскольку чрезмерная фильтрация может исказить тонкие линии связей, нарушить форму колец или привести к слиянию близко расположенных элементов. По той же причине бинаризация дополняется проверкой целостности структуры сразу после её выполнения: это позволяет отбрасывать или ослаблять такие преобразования, которые ухудшают связность и геометрию химической схемы.

Морфологические операции очистки и соединения также сохраняются. Их назначение состоит в удалении мелких артефактов и восстановлении незначительных разрывов, не разрушая при этом исходную структуру молекулы. В отличие от текстового OCR, где объединение соседних фрагментов часто полезно для восстановления символов, при обработке химических изображений слишком агрессивное склеивание может привести к ложному объединению связей, атомных обозначений и служебных элементов. Поэтому в режиме OCSR морфологические преобразования ограничиваются по силе и подчиняются требованию сохранения исходной геометрии.

Кроме того, в модуле сохраняются этапы коррекции наклона, удаления теней, удаления краевых областей, инверсии фона и добавления внутренних отступов, однако их применение рассматривается как вспомогательное и зависит от состояния конкретного изображения. В совокупности это позволяет использовать единую архитектуру адаптивной предобработки как для OCR текста, так и для OCSR, но с различной параметризацией. Таким образом, для химических структур используется тот же базовый подход, что и для текстовых изображений, однако его настройки смещены от агрессивного улучшения читаемости к умеренному сохранению графической структуры объекта.